%----------
%	ANALISIS GENERAL - MACHINE LEARNING
%----------	

A total of 255 experiments have been conducted for \textit{"Análisis General}" in the span of five months (12/02/24 - 08/07/24), defining eleven different windows that can be seen in Figures ~\ref{fig:exp_analisis_general} and ~\ref{fig:exp_analisis_general_classes}. In particular, Figure \ref{fig:exp_analisis_general} shows the overall trend in weighted average F1 scores across different experimental windows, marked by significant changes. The linear trend line, with $R^2 = 0.36$, underscores the moderate yet positive impact of these incremental improvements. Figure \ref{fig:exp_analisis_general_classes}, on the other hand, provides a class-specific breakdown, highlighting the individual performances for 'Comentario Positivo' and 'Comentario Negativo', where we can see that overall the metrics for F1-Score is lower for the latter one. 

These insights can also be seen in Table~\ref{tab:window-metrics-ageneral}, where we notice constant improvement in almost every window, having W11 as the window with best mean performance metrics, with weighted average F1-Score of 0.919 and 0.936 and 0.857 for 'Comentario Positvo' and 'Comentario Negativo', respectively.

\begin{figure}[H]
	\ffigbox[\FBwidth]
	{\caption{Distribution of the experiments conducted for "Análisis General" over time with the different windows and trend line.}\label{fig:exp_analisis_general}}
 % [Aquí ponemos como queremos que aparezca en el índice, sin la referencia]{Aquí como queremos que aparezca en el documento, con la referencia}
	{\includegraphics[scale=0.25]{Plantilla_TFG_ingles_2019/imagenes/Experiments in _Análisis General_ over time.png}}
\end{figure}

\begin{figure}[H]
	\ffigbox[\FBwidth]
	{\caption{Distribution of the experiments conducted for "Análisis General" per class over time with the different windows}\label{fig:exp_analisis_general_classes}}
 % [Aquí ponemos como queremos que aparezca en el índice, sin la referencia]{Aquí como queremos que aparezca en el documento, con la referencia}
	{\includegraphics[scale=0.25]{Plantilla_TFG_ingles_2019/imagenes/Experiments in _Análisis General_ over time_classes.png}}
\end{figure}

\begin{figure}[H]
	\ffigbox[\FBwidth]
	{\caption{Distribution of the experiments conducted for "Análisis General" per class over time with the different windows and trend line}\label{fig:exp_analisis_general_classes}}
 % [Aquí ponemos como queremos que aparezca en el índice, sin la referencia]{Aquí como queremos que aparezca en el documento, con la referencia}
	{\includegraphics[scale=0.25]{Plantilla_TFG_ingles_2019/imagenes/Experiments in _Análisis General_ over time_classes_trendline.png}}
\end{figure}


\begin{table}[H]
    \centering
    \begingroup
    \small % Change to desired font size
    \ttabbox[\FBwidth]{
        \caption{Average Metrics for Each Window in ML experiments for AG}
        \label{tab:window-metrics-ageneral} % Unique label for the table
    }{
        \begin{tabularx}{\textwidth}{
            X
            >{\centering\arraybackslash}X
            >{\centering\arraybackslash}X
            >{\centering\arraybackslash}X
            >{\centering\arraybackslash}X
            >{\centering\arraybackslash}X
            >{\centering\arraybackslash}X
            >{\columncolor{gray!15}\centering\arraybackslash}X % Highlight this column
        }
            \hline
            \rowcolor{gray!30} % Adding a gray color to the header row
            \textbf{} & \textbf{Acc. (Std.)} & \textbf{Time (s)} & \textbf{W. Avg. Precision} & \textbf{W. Avg. Recall} & \textbf{W. Avg. F1-Score} & \textbf{C. Pos. F1-Score} & \textbf{C. Neg. F1-Score} \\
            \hline
            W1 & 0.667 (0.031) & 243.190 & 0.673 & 0.667 & 0.667 & 0.715 & 0.535 \\
            \hline
            W2 & 0.534 (0.058) & 343.253 & 0.547 & 0.534 & 0.522 & 0.563 & 0.603 \\
            \hline
            W3 & 0.793 (0.008) & 3287.552 & 0.818 & 0.793 & 0.802 & 0.875 & 0.620 \\
            \hline
            W4 & 0.834 (0.019) & 1094.649 & 0.832 & 0.834 & 0.832 & 0.904 & 0.592 \\
            \hline
            W5 & 0.713 (0.038) & 479.563 & 0.806 & 0.713 & 0.718 & 0.778 & 0.428 \\
            \hline
            W6 & 0.818 (0.028) & 1065.093 & 0.825 & 0.818 & 0.817 & 0.878 & 0.601 \\
            \hline
            W7 & 0.790 (0.032) & 473.784 & 0.794 & 0.790 & 0.789 & 0.784 & 0.795 \\
            \hline
            W8 & 0.818 (0.025) & 1024.881 & 0.825 & 0.818 & 0.817 & 0.877 & 0.603 \\
            \hline
            W9 & 0.820 (0.020) & 785.324 & 0.832 & 0.820 & 0.815 & 0.879 & 0.590 \\
            \hline
            W10 & 0.798 (0.024) & 800.845 & 0.829 & 0.798 & 0.800 & 0.835 & 0.676 \\
            \hline
            W11 & 0.919 (0.012) & 5390.420 & 0.922 & 0.919 & 0.919 & 0.936 & 0.857 \\
            \hline
            \multicolumn{8}{l}{Source: 'Training.xlsx'} \\
            \hline
        \end{tabularx}
    }
    \endgroup
\end{table}





\begin{itemize}
    \item \textbf{W1: Baseline Establishment (12-02-2024):} Implementation of two fundamental models, Random Forest and Logistic Regression, to establish a baseline for performance without any modifications to the dataset.

    \item \textbf{W2-W3: Introduction of Data Balancing:}
        \begin{itemize}
            \item \textbf{W2: Down-sampling (27-02-2024):} Introduced down-sampling to handle class imbalances. This method improved the performance metrics for the negative class, as seen in Figure~\ref{fig:exp_analisis_general_classes}, enabling the models to more effectively learn and differentiate between both classes as detailed in Table~\ref{tab:window-metrics-ageneral}.
            
            \item \textbf{W3: Up-sampling (13-03-2024):} Up-sampling was introduced, which increased the weighted average F1 score. However, the improvement was primarily due to misclassifying the negative class, as highlighted in Figure~\ref{fig:exp_analisis_general_classes} and Table~\ref{tab:window-metrics-ageneral}.
        \end{itemize}

    \item \textbf{W4-W5: Introduction of Advanced Embeddings and Preprocessing:}
        \begin{itemize}
            \item \textbf{W4: (31-03-2024):} Implemented the 'Roberta' embedding described in Section~\ref{text_encoding_dl}, and included additional stylistic features detailed in Section~\ref{stylistic_features} like filtering tweets in 'es' language, adding the tweet length, and the number of adjectives to enrich the feature set. These modifications improved the overall F1 score but did not enhance metrics for 'Comentario Negativo'. Added the 'Beto' embedding to the pipeline and removed the 'Comentario Neutro' category to focus strictly on positive and negative comments. Additional data balancing techniques, such as SMOTE and ADASYN, were applied. This period saw the best mean performance metrics, achieving a weighted average F1 score of 0.889 (0.929 for 'Comentario Positivo' and 0.694 for 'Comentario Negativo').
        
            \item \textbf{W5: (13-04-2024):} Implemented the "Robertuito" embedding described in Section~\ref{text_encoding_dl}. This modification led to the worst result for 'Comentario Negativo' of the experimentation phase as seen in Table~\ref{tab:window-metrics-ageneral}, indicating that not all advanced embeddings contribute positively to the model performance.
        \end{itemize}

    \item \textbf{W6-W11: Text Pre-Processing and Final Evaluations:}
        \begin{itemize}
            \item \textbf{W6: (24-04-2024):} Applied comprehensive text preprocessing as explained in Section~\ref{sec:data_preprocessing}, including normalizing text to lowercase, removing emojis, numbers, and punctuation, eliminating non-contributive words, and reducing elongated words and repeated characters. The performance analysis remained similar to W4.
        
            \item \textbf{W7: (25-04-2024):} Utilized a focused downsampling approach and introduced the lexicons detailed in Section~\ref{lexicons}. As shown in Table~\ref{tab:window-metrics-ageneral}, the performance metrics for both classes are good and very similar, indicating that the models are effectively learning and capturing both classes.
        
            \item \textbf{W8: (26-04-2024):} Conducted a near-comprehensive test involving all data balancing methods to evaluate the effects of the previous text processing modifications. However, the results did not show improvement, as the performance for 'Comentario Negativo' decreased. Moreover, we could see in Table~\ref{tab:window-metrics-ageneral} a major increase in the mean execution time in this window.
        
            \item \textbf{W9: (03-05-2024):} Concluded the experimentation phase with no balancing techniques, serving as a control setup to evaluate the effectiveness of the changes made throughout the different windows. As expected, the absence of data balancing led to worse model performance, as seen in Table~\ref{tab:window-metrics-ageneral}.
        
            \item \textbf{W10: (08-05-2024):} Applied traditional embeddings with all balancing techniques. This phase aimed to consolidate previous findings and ensure robustness across various embeddings and balancing methods. Results were very similar as in W9.
        
            \item \textbf{W11: (03-07-2024):} Introduced OpenAI embedding. From Table~\ref{tab:window-metrics-ageneral}, the metrics and F1 scores showed better performance for both classes, demonstrating the superior quality of this embedding in capturing the particular context of the dataset, highlighting the major increase in F1-Scores of 0.936 and 0.857 for 'Comentario Positivo' and 'Comentario Negativo', respectively.
        \end{itemize}
\end{itemize}



%\color{blue}
%In Appendix \ref{appendix_results_ml_ageneral}, we provide more graphs that help understand the process. In particular, we provide class-specific plots where each subplot represents a specific model, embedding and balancing technique, providing insights into how each type affects the classification performance:

%\begin{itemize}
%\item \textbf{Figure \ref{fig:exp_analisis_general_classes_models}} displays the F1 scores across different models. The variations in the plots reflect the strengths and limitations of each model in handling the dataset and capturing the nuances of 'Comentario Positivo' and 'Comentario Negativo'.

%\item \textbf{Figure \ref{fig:exp_analisis_general_classes_balance}} illustrates the effects of different balancing techniques on the F1 scores. The results underscore the importance of balancing techniques in achieving better performance, with down-sampling and SMOTE showing notable improvements in capturing both classes accurately.

%\item \textbf{Figure \ref{fig:exp_analisis_general_classes_embeddings}} shows the F1 scores for 'Comentario Positivo' and 'Comentario Negativo' across different embedding types. The trends highlight the impact of advanced embeddings like 'Roberta' and 'Beto', and the introduction of OpenAI embeddings which significantly improved performance metrics for both classes.

%\end{itemize}

%\color{black}

Taking everything into account, we provide the following detailed tables with the best overall models for \textit{Análisis General}. Table \ref{tab:ageneral-bestmodels} lists the models, their configurations, and Table \ref{tab:ageneral-bestmodels_metrics} presents the performance metrics for these models. Notice that we have included two näive models in both tables in order to compare them to our models. {\agmajor} is a classifier that only predicts the majority class, while {\agrandom} uses the probabilities of each class to predict at random. It is important to point out that for these models we are assuming that the Standard Deviation is zero because we are using the probabilities of each class, and that F1-Score is approximately equal to such probabilities.

\begin{table}[H]
    \centering
    \begingroup
    \small % Change to desired font size
    \ttabbox[\FBwidth]{
        \caption{Best Models for AG}
        \label{tab:ageneral-bestmodels} % Unique label for the table
    }{
        \begin{tabularx}{\textwidth}{
            X  % Name
            X  % Model
            X  % Embedding
            >{\hsize=0.5\hsize\centering\arraybackslash}X  % Embed. Size
            >{\hsize=0.3\hsize\centering\arraybackslash}X  % CV
            X  % Balance
            }
            \hline
            \rowcolor{gray!30} % Adding a gray color to the header row
            \textbf{} & \textbf{Model} & \textbf{Embedding} & \textbf{Embed. Size} & \textbf{CV} & \textbf{Balance} \\
            \hline
            \multicolumn{6}{l}{\textbf{Näive Models}} \\
            \hline 
            \agmajor & - & - & - & - & - \\
            \hline
            \agrandom & - & - & - & - & - \\
            \hline
            \multicolumn{6}{l}{\textbf{Finetuned Models}} \\
            \hline
            \agrfted & Random Forest & text-embedding-3-large & - & 5 & Down-sampling \\
            \hline
            \agmlpted & MLP & text-embedding-3-large & - & 5 & Down-sampling \\
            \hline
            \agmlpteu & MLP & text-embedding-3-large & - & 5 & Up-sampling \\
            \hline
            \agmlpxlmd & MLP & XLM-RoBERTa-base & 500 & 5 & Down-sampling \\
            \hline
            \agxgbrd & XGBoost & RoBERTa & 500 & 5 & Down-sampling \\
            \hline
            \aglrbcd & Logistic Regression & Beto-cased & 500 & 5 & Down-sampling \\
            \hline
            \agxgbbma & XGBoost & BERT-multi & 500 & 5 & ADASYN \\
            \hline
            \multicolumn{6}{l}{Source: 'Training.xlsx'} \\
            \hline
        \end{tabularx}
    }
    \endgroup
\end{table}




\begin{table}[H]
    \centering
    \begingroup
    \small % Change to desired font size
    \ttabbox[\FBwidth]{
        \caption{Metrics for the Best Models for AG}
        \label{tab:ageneral-bestmodels_metrics} % Unique label for the table
    }{
        \begin{tabularx}{\textwidth}{
            X
            >{\centering\arraybackslash}X
            >{\centering\arraybackslash}X
            >{\centering\arraybackslash}X
            >{\columncolor{gray!15}\centering\arraybackslash}X % Highlight this column
        }
            \hline
            \rowcolor{gray!30} % Adding a gray color to the header row
            \textbf{} & \textbf{Acc. (Std.)} & \textbf{Weighted Avg. F1-Score} & \textbf{C. Positivo F1-Score} & \textbf{C. Negativo F1-Score} \\
            \hline
            \multicolumn{5}{l}{\textbf{Näive Models}} \\
            \hline
            \agmajor & 0.777 (0.000) & 0.777 & 1.000 & 0.000 \\
            \hline
            \agrandom & 0.641 (0.000) & 0.641 & 0.778 & 0.222 \\
            \hline
            \multicolumn{5}{l}{\textbf{Finetuned Models}} \\
            \hline
            \agrfted & 0.905 (0.031) & 0.904 & 0.899 & 0.910 \\
            \hline
            \agmlpted & 0.898 (0.019) & 0.898 & 0.894 & 0.902 \\
            \hline
            \agmlpteu & 0.945 (0.007) & 0.946 & 0.965 & 0.878 \\
            \hline
            \agmlpxlmd & 0.843 (0.027) & 0.843 & 0.839 & 0.847 \\
            \hline
            \agxgbrd & 0.836 (0.015) & 0.836 & 0.830 & 0.842 \\
            \hline
            \aglrbcd & 0.843 (0.027) & 0.843 & 0.845 & 0.841 \\
            \hline
            \agxgbbma & 0.875 (0.025) & 0.873 & 0.921 & 0.705 \\
            \hline
            \multicolumn{5}{l}{Source: 'Training.xlsx'} \\
            \hline
        \end{tabularx}
    }
    \endgroup
\end{table}




We have chosen model \textbf{\agrfted}  as the best performing model because it achieved the highest F1-Score for negative comments (Comentario Negativo F1-Score) while maintaining strong performance across other metrics. Ensuring a high Comentario Negativo F1-Score is crucial, as the second classification task, \textit{Contenido Negativo}, focuses exclusively on identifying negative comments. Although, it achieves the highest standard deviation of the models detailed in Table~\ref{tab:ageneral-bestmodels_metrics}, we can see significant improvements compared to näive models {\agmajor} and {\agrandom}.

In Figure \ref{fig:exp_analisis_general_metrics} we can see the confussion matrix for such model. We notice a higher number of true negatives (71) compared to false negatives (12), with a recall of 0.973 for "Comentario Negativo" and 0.838 for "Comentario Positivo", which implies the model is particularly effective at identifying "Comentario Negativo" instances correctly, potentially at the expense of some positive comments.


\begin{figure}[H]
	\ffigbox[\FBwidth]
	{\caption{Confusion Matrix for the best overall model (\textbf{\agrfted}) of the Machine Learning experiments conducted for "Análisis General"}\label{fig:exp_analisis_general_metrics}}
 % [Aquí ponemos como queremos que aparezca en el índice, sin la referencia]{Aquí como queremos que aparezca en el documento, con la referencia}
	{\includegraphics[scale=0.65]{Plantilla_TFG_ingles_2019/imagenes/Experiments in _Análisis General_ best_model.png}}
\end{figure}



